\section{The complete resolution profile of positive affine acquisition}
\label{sec:v18-affine}
The rank classification does not prescribe how acquired directions
vanish. For a general affine acquisition, that question has an exact
answer in terms of a cross-covariance operator. Unlike a future-test
spectrum alone, this operator contains both acquisition and observation.
No assumption on the dimension, ordering, or density of the latent
space is required.

\subsection{The experiment-level invariant}
Let $b,q\ge1$, let $f=(f_1,\ldots,f_b)^t\in C(X)^b$, and suppose
\begin{equation}\label{eq:v18-affine-bound}
 \sup_{x\in X}\sum_{i=1}^b|f_i(x)|\le a<1.
\end{equation}
The acquisition command is uniform on $U=[-1,1]^b$. The two report
likelihoods are
\begin{equation}\label{eq:v18-affine-likelihood}
 L_y(u,x)=\tfrac12(1+y\,u^tf(x)),\qquad y\in\{-1,1\}.
\end{equation}
Both are at least $(1-a)/2$. Let $H_1,\ldots,H_q\in C(X)$ be the
success probabilities of any fixed finite list of subsequent physical
queries, with $0\le H_j\le1$, and let $w_j>0$, $\sum_jw_j=1$.
The chosen query is revealed independently after the label is stored.
Put
\begin{equation}\label{eq:v18-covariance}
 h=(\sqrt{w_j}H_j)_{j=1}^q,\quad \bar h=\mu h,\quad m=\mu f,
 \qquad C_\mu=\mu[(h-\bar h)(f-m)^t]:\R^b\longrightarrow\R^q.
\end{equation}
The command cube fixes the units of $f$. We do not whiten either side
of this matrix by a possibly singular covariance. Such a whitening
would change the physical scale being classified.

Let $d=\min(b,q)$, let $s_1\ge\cdots\ge s_d\ge0$ be the singular
values of $C_\mu$, and set
\begin{equation}\label{eq:v18-affine-profile}
 V_\ell(C_\mu)=\prod_{i=1}^\ell s_i
             =\|\mathop{\bigwedge}\nolimits^\ell C_\mu\|_{\mathrm{op}},
 \qquad
 E_{C_\mu}(M)=\max_{1\le\ell\le d}
           \left(\frac{V_\ell(C_\mu)}{M}\right)^{2/\ell}.
\end{equation}
Zero products remain zero. Write $R^{\mathrm{av}}_M$ and
$R^{\mathrm{max}}_M$ for the unconditional and worst-history checkpoint
risks, with the same finite-label and randomization conventions as above.

\begin{theorem}[Uniform affine covariance classification]
\label{thm:v18-affine-classification}
There are $0<c\le C<\infty$, depending only on $a,b,q$, such that every
experiment \eqref{eq:v18-affine-likelihood}, every prior $\mu$, every
query list as above, and every integer $M\ge1$ satisfy
\begin{equation}\label{eq:v18-affine-law}
 cE_{C_\mu}(M)\le R^{\mathrm{av}}_M
 \le R^{\mathrm{max}}_M\le CE_{C_\mu}(M).
\end{equation}
The upper code can use only reachable prediction vectors. The constants
are independent of the nonzero singular values, their multiplicities,
the prior, and the query weights. In particular, the assertion is
uniform through every rank loss in $C_\mu$. If $C_\mu=0$, a single
label is exact for the prescribed query task.
\end{theorem}

\subsection{Exact acquisition geometry, including evidence}
\begin{lemma}[Projective acquisition coordinates]\label{lem:v18-projective-law}
Set $v=yu$ and $\theta=v/(1+m^tv)$. The acquired prediction vector is
exactly
\begin{equation}\label{eq:v18-projective-prediction}
 p(u,y)=\bar h+C_\mu\theta.
\end{equation}
The law of $v$ has density $(1+m^tv)/2^b$ on $U$. The image
$\Theta=\{v/(1+m^tv):v\in U\}$ satisfies
\begin{equation}\label{eq:v18-projective-sandwich}
 rB_b\subset\Theta\subset RB_b,
 \qquad r=\frac1{2(1+a)},\quad R=\frac{\sqrt b}{1-a},
\end{equation}
where $B_b$ is the Euclidean unit ball. On $\Theta$ the density of
$\theta$ is
\begin{equation}\label{eq:v18-projective-density}
 \frac{(1-m^t\theta)^{-(b+2)}}{2^b},
\end{equation}
and lies between $(1-a)^{b+2}/2^b$ and $(1+a)^{b+2}/2^b$.
\end{lemma}
\begin{proof}
Bayes' formula and the definition of covariance give
\[
 p(u,y)=\frac{\bar h+\mu(hf^t)v}{1+m^tv}
       =\bar h+\frac{C_\mu v}{1+m^tv}.
\]
For each sign $y$, the change $u\mapsto v=yu$ preserves volume and
$U$. The joint density of this branch is
$(1+m^tv)/(2\,2^b)$. Adding the two branches gives the stated density
of $v$. This calculation includes the probability of the observed
report; it is not a distribution obtained by conditioning away evidence.

Write $F(v)=v/(1+m^tv)$. Since $|m^tv|\le a$ on $U$, its inverse is
$F^{-1}(\theta)=\theta/(1-m^t\theta)$ on its image. Direct
differentiation and the rank-one determinant identity give
\[
 DF(v)=\frac{I}{1+m^tv}-\frac{vm^t}{(1+m^tv)^2},\qquad
 \det DF(v)=(1+m^tv)^{-(b+1)}.
\]
Multiplying the density of $v$ by the inverse Jacobian gives
\eqref{eq:v18-projective-density}. Its bounds follow from
$1-a\le1+m^tv\le1+a$.

The outer inclusion in \eqref{eq:v18-projective-sandwich} follows from
$\|v\|\le\sqrt b$. For the inner inclusion, $\|m\|\le\|m\|_1\le a$.
If $\|\theta\|\le r$, then $|m^t\theta|<1/2$ and
$\|\theta/(1-m^t\theta)\|\le2r\le1$. This inverse point belongs
to $B_b\subset U$, and substitution verifies that its image is
$\theta$. Thus the full ball is contained in $\Theta$.
\end{proof}

\begin{lemma}[Quantization of a linearly observed thick set]
\label{lem:v18-thick-linear}
Let $A:\R^b\to\R^q$ be linear. Suppose a random vector $\theta$ is
supported in $RB_b$, and its law dominates $\gamma$ times Lebesgue
measure on $rB_b$, where $r,R,\gamma>0$. Let $s_1\ge\cdots\ge s_d\ge0$
be the singular values of $A$, $d=\min(b,q)$. For all $M\ge1$, the
optimal mean squared $M$-center error is bounded below by a constant
times
\[
 \max_{1\le\ell\le d}(s_1\cdots s_\ell/M)^{2/\ell}.
\]
The whole image of the support has an $M$-center covering of squared
radius bounded above by another constant times the same quantity.
If that image is compact, its centers can be chosen in the image.
The constants depend only on $b,q,r,R,\gamma$ and not on $A$.
\end{lemma}
\begin{proof}
Choose orthogonal singular coordinates. Fix an index $\ell$ with
$s_\ell>0$. Project the first $\ell$ input coordinates. For every
$z\in(r/2)B_\ell$, the fiber of $rB_b$ has
$(b-\ell)$-dimensional volume at least
\[
 \omega_{b-\ell}(\sqrt3r/2)^{b-\ell},
\]
where the zero-dimensional volume is one. Consequently the law of
the first $\ell$ coordinates dominates a positive constant times
Lebesgue measure on $(r/2)B_\ell$. After multiplication by
$\operatorname{diag}(s_1,\ldots,s_\ell)$ it dominates a fixed positive
mass times the uniform measure of an ellipsoid of volume
$\omega_\ell(r/2)^\ell s_1\cdots s_\ell$. The union of $M$
projected radius-$t$ balls has volume at most $M\omega_\ell t^\ell$.
Taking $t=(r/2)(s_1\cdots s_\ell/(2M))^{1/\ell}$ leaves at least
half that uniform mass outside. This proves the lower bound for the
chosen index and hence for their maximum. It remains valid when $t$
exceeds a short ellipsoid axis: the ball-volume estimate is an upper
bound on the volume of its intersection with the ellipsoid, not an
assumption that the ball fits inside. The argument also allows arbitrary
centers and randomized encoders and decoders by conditional averaging.

For the upper bound, the image lies in the rectangle with active
half-sides $Rs_i$. A grid of mesh $\varepsilon/\sqrt d$ on each active
axis covers that rectangle by at most
\[
 \prod_{i=1}^d(1+B_ds_i/\varepsilon)
 \le A_d\sum_{j=0}^d\varepsilon^{-j}s_1\cdots s_j
\]
balls of radius $\varepsilon$, after increasing constants depending
on $d,R$. The inequality follows by expanding the product: each
$j$-fold product is at most the product of the largest $j$ singular
values. Zero axes need a single coordinate and cause no division.
Let $e=\max_j(s_1\cdots s_j/M)^{1/j}$. If $e>0$, then
$s_1\cdots s_j\le Me^j$. Choose $L$ large enough that
$A_d\sum_{j=1}^dL^{-j}\le1/2$. For $M\ge2A_d$, radius $Le$
therefore suffices with at most $M$ centers. Empty balls are discarded;
replacing a remaining center by an image point in that ball at most
doubles the radius. For $M<2A_d$, one image point has covering radius
at most $2Rs_1$, while $e\ge s_1/M$. Increasing the constant handles
these budgets. If $A=0$, the image is a singleton and every assertion
has zero error. No rank-dependent lower bound on a singular value
entered either argument.
\end{proof}

\begin{proof}[Proof of Theorem~\ref{thm:v18-affine-classification}]
Apply Lemma~\ref{lem:v18-thick-linear} to $A=C_\mu$ and the exact law
in Lemma~\ref{lem:v18-projective-law}. The constants $r,R,\gamma$ in
that application depend only on $a,b$. Translation by $\bar h$ changes
neither coding nor error. Weighted query probabilities identify physical
squared-loss excess with the Euclidean norm used in those lemmas, so no
inverse query weight is paid in the constants. Every representative
belongs to $\bar h+C_\mu\Theta$ and hence is a posterior prediction
vector of an admitted history. Its unweighted entries are valid query
probabilities. This proves the upper bound for every history, while
the exact unconditional law proves the lower bound. If $C_\mu=0$,
\eqref{eq:v18-projective-prediction} is independent of the history and
its known constant prediction requires one label.
\end{proof}

\subsection{Operational completeness and an interior degeneration}
\begin{corollary}[Complete invariant in the affine class]
\label{cor:v18-affine-invariant}
In any family with fixed bounds on $a<1$, $b$, and $q$, the products
$V_\ell(C_\mu)$ characterize uniform comparison of entire integer-budget
checkpoint-risk curves. More precisely, after padding to a fixed maximum
dimension, two such families have uniformly comparable curves if and
only if their corresponding products are uniformly comparable,
including matching zero products. Moreover
\begin{equation}\label{eq:v18-affine-inverse}
 \inf_{M\in\N,\ M\ge1}M\bigl(R^{\mathrm{av}}_M\bigr)^{\ell/2}
 \asymp V_\ell(C_\mu).
\end{equation}
For positive rank $r_\mu$, the eventual logarithmic slope is $-2/r_\mu$.
\end{corollary}
\begin{proof}
The singular values are ordered and Theorem~\ref{thm:v18-affine-classification}
provides the forward representation uniformly for this entire experiment
class. Lemma~\ref{lem:integer-envelope-duality} gives, for
$e(M)=E_{C_\mu}(M)^{1/2}$,
\[
 V_\ell\le\inf_{M\ge1}Me(M)^\ell\le2V_\ell.
\]
Combining it with \eqref{eq:v18-affine-law} proves
\eqref{eq:v18-affine-inverse}. Comparison of risks therefore implies
comparison of products; conversely compare each term in
\eqref{eq:v18-affine-profile} and take their finite maximum.
The argument includes zero products by the same envelope identity.
For fixed positive rank, the last positive branch dominates for all
sufficiently large budgets and yields the slope. The all-zero case is
handled by the one-label assertion of the theorem, not by taking a
logarithm of zero.
\end{proof}

This is an experiment-level classification: the covariance is computed
from likelihood features and physical future tests before the coding
problem is solved. It does not posit an ordered-envelope law and then
use that law to discover the forward geometry. The inverse still uses
the full unbounded sequence of integer budgets; it is not reconstruction
of an unknown experiment from a finite sample of observed errors.

\begin{proposition}[Rank change among full-support priors]
\label{prop:v18-interior-example}
Let $X=[-1,1]$, use $f(x)=\alpha x$ with $0<\alpha<1$, and let the
future query be a Bernoulli experiment with success probability
$H(x)=1/2+\beta x^2$, where $0<\beta<1/2$. For the full-support priors
\[
 d\mu_t(x)=(1+tx)\,dx/2,\qquad |t|\le t_0<1,
\]
the sharp acquisition-checkpoint law is
\begin{equation}\label{eq:v18-interior-law}
 R^{\mathrm{av}}_M\asymp R^{\mathrm{max}}_M
 \asymp (\alpha\beta)^2t^2M^{-2},\qquad M\ge1,
\end{equation}
uniformly in $t$, including zero. At $t=0$ one label is exact for the
specified future query; at every $t\ne0$ the attainable rank is one.
\end{proposition}
\begin{proof}
Direct integration gives
\[
 \mu_tx=t/3,\qquad \mu_tx^2=1/3,\qquad \mu_tx^3=t/5,
 \qquad \Cov_{\mu_t}(x^2,x)=4t/45.
\]
Thus the sole singular value of \eqref{eq:v18-covariance} is
$4\alpha\beta|t|/45$. Theorem~\ref{thm:v18-affine-classification}
proves \eqref{eq:v18-interior-law}. Equivalently, with $v=yu$, the
posterior query probability is exactly
\[
 \frac12+\beta\,
 \frac{1/3+\alpha vt/5}{1+\alpha vt/3},
\]
whose derivative in $v$ is
$4\alpha\beta t/[45(1+\alpha vt/3)^2]$.
It is zero throughout the command interval precisely at $t=0$.
Both acquisition reports and both future outcomes have uniformly
positive likelihoods. Thus the degeneration occurs inside a dominated
full-support prior family; it is not loss of support or vanishing
evidence. The centered moment equation $\mu x^3-\mu x\,\mu x^2=0$
is exactly its rank stratum.
\end{proof}

The example exhibits a degeneration governed by the interaction of
acquisition and observation rather than a collision of future exponents.
Its role is to make the structural invariant concrete; the uniform
classification applies to arbitrary dimensions and arbitrary latent
spaces in \eqref{eq:v18-affine-bound}, not only to this example. The
multi-step collision and contrast laws are retained below and are not
replaced by a one-step assertion.
